% generated by intro-ledger; compile inside your paper with \bibliography{references}
\section{1 Introduction}

Rapid landslide mapping after a triggering event informs emergency response, and two kinds of imagery compete for the job. Synthetic Aperture Radar (SAR) can provide measurements of affected areas independent of weather or lighting conditions \cite{boehm2022}, yet optical images gave an overall accuracy of 99.20\% in a Hokkaido earthquake case where SAR-based CNNs reached overall accuracies beyond 94\% \cite{nava2021}. Choosing a sensor therefore trades accuracy against availability, and a model that serves both would shorten the time to a usable inventory.

\section{2 Related Work}

**Input modality.** Optical and SAR inputs have been compared directly: optical CNNs were more accurate, whereas SAR still allowed rapid mapping during storms and under dense cloud cover \cite{nava2021}. Boehm et al. show that deep learning models can detect landslides from SAR data, with an area under the precision-recall curve above 0.7, and that early detection is possible when SAR is combined with terrain information \cite{boehm2022}. Combining Sentinel-1 and Sentinel-2 bands gave the highest F1-score (0.69) in a heterogeneous ensemble, compared with 0.61 for Sentinel-2 bands only \cite{ganerd2023}.

**Benchmarks and architecture comparisons.** Landslide4Sense provides 3,799 image patches that fuse Sentinel-2 layers with elevation and slope layers, and the authors found that ResU-Net outperformed ten other segmentation models \cite{ghorbanzadeh2022}. LMHLD is larger, with 25,365 patches from five sensors across seven study areas \cite{liu2025}. Oak et al. instead compare four CNN-based models on one dataset, where LinkNet gave the best F1-score of 85.7\% \cite{oak2024}. These resources differ in scale, but each compares models within a given input configuration rather than across sensors.

**Input design.** Ahmad et al. identify a compact 8-channel subset that matches or exceeds the F1 of configurations using up to 30 channels \cite{ahmad2026}, while Chen and Fan study how selecting contributing factors changes the accuracy of machine learning and deep learning susceptibility models \cite{chen2023}. Both treat input selection within a fixed data setting.

\section{3 Research gap and contributions}

Three limitations stand out. Using SAR is hindered by the domain knowledge needed for pre-processing and interpretation \cite{boehm2022}; deep learning landslide mapping models often show sensitivity problems, overfitting and low accuracy \cite{ganerd2023}; and the effect of factor selection on model accuracy is still unclear \cite{chen2023}. We argue that a compact model shared across optical and SAR input, adaptable to a new region with few labelled tiles, is missing.

Our contributions are: (1) LiteSlide, a transfer model with a shared encoder and 4.1 million parameters; (2) a mean F1 of 0.812 on four held-out regions, against 0.771 for the best of three baselines; and (3) an SAR-only F1 of 0.703, and a drop of 0.04 in F1 when fine-tuning tiles fall from 200 to 50.
